\documentclass[11pt]{article}
\usepackage[T1]{fontenc}
\usepackage{textcomp}
\usepackage[margin=0.75in]{geometry}
\usepackage{amsmath,amssymb,amsthm}
\usepackage{array,booktabs}
\usepackage{hyperref}
\setlength{\parindent}{0pt}
\newtheorem{theorem}{Theorem}
\theoremstyle{definition}
\newtheorem{definition}{Definition}
\title{Flow Proof Artifact: Group-inverse-agreement}
\author{Generated by \texttt{flow doc proof}}
\date{Flow Proof Book}
\begin{document}
\maketitle
Left and right inverses coincide in a group.\\[0.5em]
\textbf{Source.} dummit-foote — *Abstract Algebra*, §1.1\\[0.5em]
\bigskip
\noindent{\large \textbf{Derived fact 1.} \textit{If h * g = 1 and g * inv(g) = 1 then h = inv(g)} \hfill Group \cdot inverse \cdot left and right inverses agree}\par
\smallskip\noindent\textit{Coordinate: Group \cdot inverse \cdot left and right inverses agree}\par
\smallskip\noindent\textit{Source: dummit-foote}\par
\smallskip\noindent\textit{Built on: inverses are unique, for inverse on Group, right inverse recovers the identity, for inverse on Group}\par
\medskip
\noindent\textbf{Goal.} If h * g = 1 and g * inv(g) = 1 then h = inv(g)\par
\noindent$\displaystyle \forall g \in Group \forall h \in Group\quad h = inv(g)$\par
\medskip
\noindent
\begin{tabular}{@{} >{\raggedright\arraybackslash}p{0.06\textwidth} >{\raggedright\arraybackslash}p{0.47\textwidth} >{\raggedleft\arraybackslash}p{0.41\textwidth} @{}}
\textbf{\#} & \textbf{Proof} & \textbf{Mathematics} \\
\midrule
\textbf{1.} & We prove that left and right inverses agree for inverse on Group. & \\[0.45em]
\textbf{2.} & We split into exhaustive cases — the claim must hold in each one. & \\[0.45em]
\textbf{3.} & Case 1 (see step 2): suppose h * g  equals  1. & \\[0.45em]
\textbf{4.} & Case 2 (see step 2): suppose g * inv(g)  equals  1. & \\[0.45em]
\textbf{5.} & We invoke the derived fact governing inverse on Group: inverses are unique, for inverse on Group (instantiated for g, h, inv(g)). & \\[0.45em]
\textbf{6.} & We invoke the derived fact governing inverse on Group: right inverse recovers the identity, for inverse on Group (instantiated for g). & \\[0.45em]
\textbf{7.} & From step 4, step 5, and step 6, this implies h equals inv(g). Together with the other cases (step 3 and step 4), the goal is discharged. Hence proven. & $\displaystyle h = inv(g)$ \\[0.45em]
\bottomrule
\end{tabular}
\label{thm:Group-inverse-left_and_right_inverses_agree}
\par\medskip\hrule
\end{document}
